% GENERATED FILE. Edit canonical lessons, then run npm run generate:books.
% canonical-source-hash: fnv1a64:5872859f937a5570
% canonical-lessons: PA-W14-timed-repair-return, PA-R159-people-quality-recall, PA-R159-strength-weakness-recall, PA-R159-hide-and-quality-mix

\chapter{Return to an Earlier Repair and Familiar Words}
\label{ch:pa-earlier-repair-and-familiar-words}
\hlchaptermodality{159}

\begin{chapteropening}
\textbf{By the end of this chapter:} \emph{I can keep a later Gurmukhi repair separate from a timed first attempt and retrieve familiar words for rich, poor, young, strong, weak, and to hide without mistaking these reviews for an A1 mock.}

From short familiar cues, recall the five earlier descriptions and the earlier verb for hiding without adding Punjabi language or borrowing scored writing credit.
\end{chapteropening}

\section[\texorpdfstring{\pa{ਨਾਂ}: \pa{ਅਮਨ} (nāṁ: aman)}{ਨਾਂ: ਅਮਨ (nāṁ: aman)}]{A later correction stays on another page}

\label{lesson:PA-W14-timed-repair-return}



\begin{quote}
Cover the earlier answer. Read the familiar label \textbf{\pa{ਨਾਂ}} and remember the A · \pa{ਕ} card value. No new card or person is being introduced.
\end{quote}

\subsection*{Your turn}
Set a 45-second \emph{practice} clock. Write the one known line in Gurmukhi, then stop when the clock ends:

\begin{quote}
\pa{ਨਾਂ} (A · \pa{ਕ}): \_\_\_\_\_\_
\end{quote}

After stopping, check \textbf{\pa{ਨਾਂ}: \pa{ਅਮਨ}}. If the value or spelling slipped, write one corrected line on a second page. Keep the first practice line and the earlier timed form visible and unchanged. Say which is the later correction; it cannot replace either first attempt.

\begin{tcolorbox}[breakable,colback=teal!4,colframe=teal!35!black,title={Before you move on}]
Cover the check. Recall the order: first attempt, stop, inspect, then repair separately. The repair helps learning but never raises the first attempt's score.
\end{tcolorbox}
\section[\texorpdfstring{\pa{ਅਮੀਰ} · \pa{ਗ਼ਰੀਬ} · \pa{ਜਵਾਨ} (amīr · garīb …)}{ਅਮੀਰ · ਗ਼ਰੀਬ · ਜਵਾਨ (amīr · garīb …)}]{Three old words for describing people}

\label{lesson:PA-R159-people-quality-recall}



\begin{quote}
Hear the cue \textbf{rich}. Say the old Punjabi word before the spoken model: \textbf{\pa{ਅਮੀਰ}} (\emph{amīr}).
\end{quote}

\subsection*{Your turn}
Say only the old describing word for each familiar meaning:

\begin{itemize}
\raggedright
  \item \textbf{poor} — \textbf{\pa{ਗ਼ਰੀਬ}} (\emph{garīb}).
  \item \textbf{young} — \textbf{\pa{ਜਵਾਨ}} (\emph{javān}).
  \item \textbf{rich} — \textbf{\pa{ਅਮੀਰ}} (\emph{amīr}).
\end{itemize}

These are word meanings, not labels for a real person in the room.

\begin{tcolorbox}[breakable,colback=teal!4,colframe=teal!35!black,title={Before you move on}]
Hear \textbf{young}, \textbf{rich}, then \textbf{poor}. Pause after each cue and say its Punjabi word before the model: \textbf{\pa{ਜਵਾਨ}}, \textbf{\pa{ਅਮੀਰ}}, \textbf{\pa{ਗ਼ਰੀਬ}}.
\end{tcolorbox}
\section[\texorpdfstring{\pa{ਤਕੜਾ} · \pa{ਕਮਜ਼ੋਰ} (takṛā · kamzor)}{ਤਕੜਾ · ਕਮਜ਼ੋਰ (takṛā · kamzor)}]{Two old words for strength}

\label{lesson:PA-R159-strength-weakness-recall}



\begin{quote}
Hear the cue \textbf{strong}. Say the old Punjabi word before the spoken model: \textbf{\pa{ਤਕੜਾ}} (\emph{takṛā}).
\end{quote}

\subsection*{Your turn}
Hear the contrast, then say the familiar word after a pause:

\begin{itemize}
\raggedright
  \item \textbf{weak} — \textbf{\pa{ਕਮਜ਼ੋਰ}} (\emph{kamzor}).
  \item \textbf{strong} — \textbf{\pa{ਤਕੜਾ}} (\emph{takṛā}).
\end{itemize}

The cues ask for meanings already learned, not a new sentence pattern.

\begin{tcolorbox}[breakable,colback=teal!4,colframe=teal!35!black,title={Before you move on}]
Hear \textbf{strong}, then \textbf{weak}. Pause after each cue and say its Punjabi word before the model: \textbf{\pa{ਤਕੜਾ}}, \textbf{\pa{ਕਮਜ਼ੋਰ}}.
\end{tcolorbox}
\section[\texorpdfstring{\pa{ਲੁਕਣਾ} (lukṇā)}{ਲੁਕਣਾ (lukṇā)}]{Hide, then mix the familiar descriptions}

\label{lesson:PA-R159-hide-and-quality-mix}



\begin{quote}
Hear the cue \textbf{to hide}. Say the earlier verb before the spoken model: \textbf{\pa{ਲੁਕਣਾ}} (\emph{lukṇā}).
\end{quote}

\subsection*{Your turn}
Mix the old meanings. Pause after each cue before its spoken model:

\begin{itemize}
\raggedright
  \item \textbf{young} — \textbf{\pa{ਜਵਾਨ}} (\emph{javān}).
  \item \textbf{to hide} — \textbf{\pa{ਲੁਕਣਾ}} (\emph{lukṇā}).
  \item \textbf{weak} — \textbf{\pa{ਕਮਜ਼ੋਰ}} (\emph{kamzor}).
  \item \textbf{rich} — \textbf{\pa{ਅਮੀਰ}} (\emph{amīr}).
  \item \textbf{strong} — \textbf{\pa{ਤਕੜਾ}} (\emph{takṛā}).
  \item \textbf{poor} — \textbf{\pa{ਗ਼ਰੀਬ}} (\emph{garīb}).
\end{itemize}

Every Punjabi word on this list was already taught. The order, not the language, is new.

\begin{tcolorbox}[breakable,colback=teal!4,colframe=teal!35!black,title={Before you move on}]
Hear \textbf{poor}, \textbf{strong}, \textbf{to hide}, \textbf{rich}, \textbf{weak}, then \textbf{young}. Pause after each cue and recall its Punjabi word before the model: \textbf{\pa{ਗ਼ਰੀਬ}}, \textbf{\pa{ਤਕੜਾ}}, \textbf{\pa{ਲੁਕਣਾ}}, \textbf{\pa{ਅਮੀਰ}}, \textbf{\pa{ਕਮਜ਼ੋਰ}}, \textbf{\pa{ਜਵਾਨ}}. Revisit only a missed old word.
\end{tcolorbox}
